% related_work.tex
%
% Subsections:
%   2.1  LLM-Guided Code Optimization
%   2.2  LLM-Guided Prompt Optimization
%   2.3  Statistical Evaluation in ML Systems
%   2.4  Evolutionary and Population-Based Search
%   2.5  Positioning: How anneal differs

\section{Related Work}
\label{sec:related-work}

\subsection{LLM-Guided Code Optimization}
\label{sec:related-code}

% AlphaEvolve (Novikov et al., 2025) — evolutionary code optimization via Gemini
% Key points: evolutionary search over code, LLM as mutation operator, no stat testing
\textbf{AlphaEvolve}~\cite{novikov2025alphaevolve} uses an LLM-guided evolutionary
algorithm to optimize algorithmic code, achieving state-of-the-art results on
mathematical programming tasks.
\TODO{Expand: what does AlphaEvolve do well? Where does it lack statistical
rigor or structured search? 3-4 sentences.}

% Other code optimization systems
\TODO{Add coverage of: FunSearch (Romera-Paredes et al., 2024), CodeEvolve if exists,
LLaMEA if applicable. 2-3 paragraphs.}

\subsection{LLM-Guided Prompt Optimization}
\label{sec:related-prompt}

% Promptbreeder (Fernando et al., 2023) — self-referential prompt evolution
\textbf{Promptbreeder}~\cite{fernando2023promptbreeder} evolves both task prompts
and mutation operators through self-referential improvement.
\TODO{Expand: self-referential aspect, no statistical significance testing,
single-objective. 3-4 sentences.}

% DSPy (Khattab et al., 2023) — declarative prompt programming
\textbf{DSPy}~\cite{khattab2023dspy} introduces a declarative programming model
for LLM pipelines, separating program structure from prompting strategy.
\TODO{Expand: DSPy teleprompters as optimization, structured but not evolutionary.
How does anneal's approach differ? 3-4 sentences.}

% EvoPrompt (Guo et al., 2023) — evolutionary prompt engineering
\textbf{EvoPrompt}~\cite{guo2023evoprompt} applies genetic algorithms and
differential evolution directly to discrete prompt optimization.
\TODO{Expand: discrete evolution, no knowledge accumulation, no stat testing.}

% OPRO (Yang et al., 2023) — optimization by prompting
\textbf{OPRO}~\cite{yang2023opro} frames optimization as a prompting task,
where the LLM itself acts as the optimizer given a natural language description
of the objective.
\TODO{Expand: natural language optimization trajectory, limitations in
multi-criteria and stochastic settings. 3-4 sentences.}

\TODO{Add: APE (Automatic Prompt Engineer), ProTeGi, TextGrad if relevant.}

\subsection{Statistical Evaluation in ML Systems}
\label{sec:related-stats}

% Wilcoxon signed-rank test in NLP evaluation
\TODO{Cite: Dror et al. (2018) "The Hitchhiker's Guide to Testing Statistical
Significance in Natural Language Processing". Discuss why non-parametric tests
are needed for bounded evaluation metrics.}

% Holm-Bonferroni correction for multiple comparisons
\TODO{Cite: Holm (1979). Discuss how multiple comparison problems arise when
evaluating a system across multiple benchmarks.}

% Bootstrap confidence intervals
\TODO{Cite: Efron and Tibshirani (1994). Discuss use in ML system evaluation.}

\subsection{Evolutionary and Population-Based Search}
\label{sec:related-evolution}

% General evolutionary algorithms background
\TODO{Brief survey: CMA-ES, NSGA-II, island models. Emphasize how anneal adapts
these ideas to discrete artifact spaces.}

% Thompson sampling for bandit problems
\TODO{Cite: Thompson (1933), Chapelle \& Li (2011). Discuss relation to
anneal's strategy selector.}

\subsection{Positioning}
\label{sec:related-positioning}

% Summary table comparing anneal to related systems across key dimensions
\begin{table}[h]
  \centering
  \caption{Comparison of LLM-guided optimization systems across key dimensions.
    \checkmark = fully supported, $\circ$ = partial, --- = not supported.}
  \label{tab:comparison}
  \TODO{Fill in table after verifying feature matrix against codebase (Task 5.1).}
  \begin{tabular}{lccccc}
    \toprule
    & \rotatebox{60}{Statistical Tests}
    & \rotatebox{60}{Structured Search}
    & \rotatebox{60}{Knowledge Accum.}
    & \rotatebox{60}{Multi-Criteria}
    & \rotatebox{60}{Cost Control} \\
    \midrule
    AlphaEvolve~\cite{novikov2025alphaevolve}   & --- & \checkmark & --- & --- & --- \\
    Promptbreeder~\cite{fernando2023promptbreeder} & --- & \checkmark & --- & --- & --- \\
    DSPy~\cite{khattab2023dspy}                 & --- & $\circ$    & --- & $\circ$ & --- \\
    EvoPrompt~\cite{guo2023evoprompt}           & --- & \checkmark & --- & --- & --- \\
    OPRO~\cite{yang2023opro}                    & --- & ---        & $\circ$ & --- & --- \\
    \textsc{anneal} (this work)                 & \checkmark & \checkmark & \checkmark & \checkmark & \checkmark \\
    \bottomrule
  \end{tabular}
\end{table}

\TODO{Write 2-3 paragraph narrative discussing what distinguishes anneal from
each category of prior work. Emphasize: (1) statistical rigor as a first-class
concern, (2) artifact-agnostic abstraction, (3) knowledge accumulation across runs.}
